\begin{table}[!tbp]
  \centering
  \caption{滞后权重比值的实证区间与各误差机制的预测（NHSN 流感流行初期，$N=1{,}267$ 个州周）。$r_1$、$r_2$ 为只用当周、三周等权两种权重相对主规格的斜率比，$q_1$、$q_2$ 为对应的截距比，$a$ 为主规格截距。实证区间为按州聚类 $B=2000$ 次重抽的 95\% 百分位区间，每次重抽都重估 $\overline{R}$。平稳环境模型（M-env）的增长期预测有两种，一种是每周一代的更新仿真（$\phi=0$，9 个情景格子，每格 400 次），另一种是 $\phi=0$ 的日分辨率更新过程。后者取流感世代间隔，环境变异系数平方 CV$^2$ 4 档、世代间隔 3 档、增长率 3 档；不含入院延迟时 36 格，加入无、短、长三档入院延迟后共 108 格，每格 300 次。乘性报告噪声模型（M-mult）的平稳行取报告冲击自相关 $\psi$ 遍历 $[0,1)$ 的解析范围（命题~\ref{prop:report-noise}(vi)）。其增长期行取 $v_R=0$ 与对数一阶自回归报告冲击，$\psi\in\{0,0.3,0.6,0.9\}$、$\tau\in\{0.15,0.3\}$、入院延迟 2 档、增长率 3 档，共 48 格，每格 300 次。其余三类为不含环境方差的二项报告模型（M-bin0）、纯测量误差模型（M-meas）与分母不含计数噪声的模型（M-nodn）。区间型预测为各格模拟平均的最小值至最大值。“--”表示该仿真未计算斜率比，“无截距”表示该机制预测 $a=0$、比值无定义。粗体为与实证区间不相交的预测。}
  \label{tab:ratios}
  \footnotesize
  \setlength{\tabcolsep}{4pt}
  \begin{tabular}{@{}lccccc@{}}
    \toprule
    机制 & $r_1$ & $r_2$ & $q_1$ & $q_2$ & $a$ \\
    \midrule
    实证 & 1.397 & 0.859 & 0.688 & 1.539 & 0.0991 \\
    \quad 95\% 百分位区间 & 1.229--1.623 & 0.599--1.083 & 0.658--0.727 & 1.420--1.656 & 0.0837--0.1165 \\
    \midrule
    M-env，平稳 & 1.419--1.463 & 0.852--0.866 & \textbf{1} & \textbf{1} & $>0$ \\
    M-env，增长期，每周一代 & 1.450--1.630 & 0.693--0.763 & \textbf{0.857--0.937} & \textbf{1.226--1.333} & $>0$ \\
    M-env，增长期，日更新无延迟 & -- & -- & \textbf{0.746--0.887} & 1.319--1.484 & $>0$ \\
    M-env，增长期，日更新全部 & -- & -- & 0.601--0.887 & 1.187--1.568 & $>0$ \\
    M-mult，平稳 & 1.462 & 0.852 & \textbf{0.883--1.338} & \textbf{0.892--1.374} & $>0$ \\
    M-mult，增长期 & -- & -- & 0.686--1.224 & 0.865--1.497 & $>0$ \\
    M-bin0 & 1.432 & 0.862 & 无截距 & 无截距 & \textbf{0} \\
    M-meas & 1.462 & 0.852 & 无截距 & 无截距 & \textbf{0} \\
    M-nodn & \textbf{1.000} & 1.000 & \textbf{1} & \textbf{1} & $>0$ \\
    \bottomrule
  \end{tabular}
\end{table}
